\chapter{The OULAD Source Tables and What Each Row Means}
\label{ch:oulad-data}

\begin{learningobjectives}
After completing this chapter, you will be able to:
\begin{itemize}
  \item identify the seven CSV tables used by the project;
  \item define module, presentation, student attempt, assessment, resource, and
    interaction in the OULAD context;
  \item state what one row means (its observation grain) and which columns
    should identify it (its candidate key) for each source table;
  \item trace important relationships among the tables;
  \item interpret relative-day fields and the question-mark missing marker; and
  \item identify limitations that must remain visible during analysis.
\end{itemize}
\end{learningobjectives}

\section{Read one CSV row before reading seven tables}

A CSV file is a text file that represents a table. Consider this simplified
excerpt from \projectfile{courses.csv}:

\begin{lstlisting}[caption={A header and one illustrative CSV data row}]
code_module,code_presentation,module_presentation_length
AAA,2013J,268
\end{lstlisting}

\begin{description}
  \item[Header] The first line names the columns. It describes what each
    position means.
  \item[Row] The second line is one data record. Here it describes one delivery
    of module AAA.
  \item[Column] Values in the same comma-separated position belong to the same
    variable. The third column records presentation length.
  \item[Cell or field] One value at the intersection of a row and column, such
    as \texttt{268}.
  \item[Identifier or ID] A value used to name or match a recorded entity. AAA
    identifies a module, but it is not a person's identity.
\end{description}

The row means that module AAA was delivered in presentation \texttt{2013J} and that the
recorded presentation length is 268 days. Reading a table starts with that
sentence, not with calculation.

\begin{validationbox}
Point to each part of the sample and answer: Which line is the header? What does
the one data row represent? Which value records presentation? Which value is a
measurement rather than an identifier?
\end{validationbox}

\section{Where the data came from: provenance}

\term{Provenance} means the recorded history of where data came from and how it
reached the project. Provenance matters because two files with the same name
can contain different bytes or versions.

OULAD was created from anonymized Open University records and released through
the UCI Machine Learning Repository \citep{oulad_dataset,oulad_paper}. The
project downloads the official archive rather than committing the source data.
An \term{archive} is one packaged file containing other files, commonly a ZIP
file. \term{Anonymized} means direct identities have been removed or replaced;
it does not mean the records have no privacy risk. The Creative Commons
Attribution 4.0 license states the conditions under which the dataset can be
reused.

The implemented downloader checks:
\begin{itemize}
  \item the archive's SHA-256 checksum;
  \item the presence of all seven CSV files and \texttt{OULAD.names}; and
  \item the exact expected archive before extraction, which means unpacking the
    files it contains.
\end{itemize}

\begin{plainlanguage}
A checksum is a fingerprint of a file. If the downloaded bytes change, the
fingerprint changes. Matching the expected checksum does not prove that every
value is correct, but it does show that the project received the exact archive
it was designed to process.
\end{plainlanguage}

\section{OULAD vocabulary}

\begin{description}
  \item[Module] An anonymized course code, represented as AAA through GGG.
  \item[Presentation] One delivery of a module in a period such as 2013J or
    2014B. The code combines a year with a start term.
  \item[Student-module-presentation attempt] One student identifier associated
    with one module and one presentation. A student can appear in more than one
    attempt.
  \item[Assessment] A TMA, CMA, or Exam defined for a module-presentation.
  \item[VLE resource] An online site or activity made available in a
    module-presentation.
  \item[VLE interaction] A recorded number of clicks by a student on a site on a
    course-relative day.
\end{description}

\section{The seven source files}

\small
\begin{longtable}{L{0.19\textwidth}L{0.13\textwidth}L{0.26\textwidth}L{0.31\textwidth}}
\caption{Verified OULAD source files used by the project}\label{tab:source-files}\\
\toprule
Source file & Rows & Observation grain & Important identifiers and purpose \\
\midrule
\endfirsthead
\toprule
Source file & Rows & Observation grain & Important identifiers and purpose \\
\midrule
\endhead
{\scriptsize\projectfile{courses.csv}} & 22 & One module-presentation & Module, presentation, presentation length \\
{\scriptsize\projectfile{assessments.csv}} & 206 & One assessment & Assessment ID plus owning module-presentation \\
{\scriptsize\projectfile{vle.csv}} & 6,364 & One site in a module-presentation & Site ID, activity type, availability weeks \\
{\scriptsize\projectfile{studentInfo.csv}} & 32,593 & One student-module-presentation attempt & Attempt context, demographics, final result \\
{\shortstack[l]{\scriptsize\texttt{student}\\[-1pt]\scriptsize\texttt{Registration.csv}}} & 32,593 & One student-module-presentation registration & Registration and unregistration relative days \\
{\scriptsize\projectfile{studentAssessment.csv}} & 173,912 & One recorded student-assessment submission & Assessment, student, submitted day, banked flag, score \\
{\scriptsize\projectfile{studentVle.csv}} & 10,655,280 & One student-site-relative-day interaction row & Module-presentation, student, site, day, click count \\
\bottomrule
\end{longtable}
\normalsize

\verified{The seven CSV files contain 10,900,970 rows in total. These counts
come from \projectfile{reports/tables/source_file_profile.csv}.}

\section{A relational map}

The seven files are \term{related}: a value in one table can be matched with a
corresponding value in another. For example, an assessment record names the
module and presentation that own it. The arrows in
\cref{fig:source-erd} mean ``can be connected through matching identifiers.''
They do not mean that every possible connection is automatically correct.

\begin{figure}[htbp]
\centering
\resizebox{\textwidth}{!}{
\begin{tikzpicture}[
  entity/.style={
    rectangle,
    draw=CourseBlue,
    fill=CourseLight,
    rounded corners=1mm,
    minimum width=34mm,
    minimum height=12mm,
    align=center,
    font=\small
  },
  link/.style={-{Latex[length=2mm]},thick,color=CourseGold}
]
\node[entity] (courses) {courses\\module-presentation};
\node[entity,right=20mm of courses] (assessments) {assessments\\assessment};
\node[entity,right=20mm of assessments] (submissions) {studentAssessment\\student-assessment};
\node[entity,below=20mm of courses] (info) {studentInfo\\student attempt};
\node[entity,right=20mm of info] (registration) {studentRegistration\\registration};
\node[entity,below=20mm of info] (studentvle) {studentVle\\interaction row};
\node[entity,right=20mm of studentvle] (vle) {vle\\presentation-site};

\draw[link] (courses) -- node[above,font=\scriptsize]{defines} (assessments);
\draw[link] (assessments) -- node[above,font=\scriptsize]{receives} (submissions);
\draw[link] (courses) -- node[left,font=\scriptsize]{contains} (info);
\draw[link] (info) -- node[above,font=\scriptsize]{matches} (registration);
\draw[link] (info) -- node[left,font=\scriptsize]{produces} (studentvle);
\draw[link] (vle) -- node[above,font=\scriptsize]{describes site} (studentvle);
\draw[link] (courses) -- node[right,font=\scriptsize]{publishes} (vle);
\end{tikzpicture}
}
\caption{Conceptual relationships among the seven source tables. Arrows show
how identifiers can connect tables. The original CSV files do not enforce
these matches as database rules.}
\label{fig:source-erd}
\end{figure}

\section{Grain before keys}

\term{Observation grain} answers: ``What does one row represent?'' A
\term{candidate key} answers: ``Which columns should uniquely identify that
row?'' A candidate key is a proposed identifier until the data is tested for
duplicates. The grain comes first. Without it, a key is merely a list of
columns.

Return to the sample row \texttt{AAA,2013J,268}. Its grain is one module
presentation. Module alone cannot identify the row because module AAA can be
delivered more than once. Module plus presentation is the candidate key.

\begin{table}[htbp]
\centering
\caption{Source-grain reasoning}
\label{tab:source-grains}
\begin{tabularx}{\textwidth}{L{0.25\textwidth}L{0.34\textwidth}Y}
\toprule
Source & Candidate identifying columns & Important caution \\
\midrule
\texttt{courses} & module + presentation & Module alone is not unique \\
\texttt{studentInfo} & module + presentation + student & Student alone is not an attempt key \\
\texttt{studentRegistration} & module + presentation + student & Must be checked against attempts \\
\texttt{assessments} & assessment ID & Module context is carried on the assessment \\
\texttt{studentAssessment} & assessment ID + student & An absent row means no recorded submission \\
\texttt{vle} & site + module + presentation & Site context belongs to a presentation \\
\texttt{studentVle} & student + site + module + presentation + day, approximately & Duplicate source-grain rows can occur; the core table deliberately has no declared primary key \\
\bottomrule
\end{tabularx}
\end{table}

\begin{validationbox}
For every table, say the following sentence aloud:

\begin{quote}
``One row represents \underline{\hspace{4cm}}, and it should be identified by
\underline{\hspace{4cm}}.''
\end{quote}

If you cannot fill both blanks precisely, do not join the table to another one
or aggregate it into summaries yet.
\end{validationbox}

\section{Time is relative}

OULAD uses days relative to the start of a presentation rather than calendar
dates. Important fields include:
\begin{itemize}
  \item \texttt{date\_registration}: registration day relative to course start;
  \item \texttt{date\_unregistration}: recorded unregistration day;
  \item \texttt{assessments.date}: expected assessment due day;
  \item \texttt{date\_submitted}: recorded submission day; and
  \item \texttt{studentVle.date}: interaction day.
\end{itemize}

Negative days are meaningful. The saved quality profile finds 688,988 VLE
activity rows before day zero. Those rows are not automatically errors: they
can represent legitimate pre-course access.

\section{Missing values and absent rows}

Two different forms of missingness matter:

\begin{enumerate}
  \item A source cell can contain the marker \texttt{?}. The staging layer
    converts it to SQL \texttt{NULL}. NULL is PostgreSQL's marker for a missing
    or unknown value; it is not the number zero or an empty string.
  \item A whole expected row can be absent. For example, a student-assessment
    pair with no matching submission row represents a missing recorded
    submission in a later analysis table.
\end{enumerate}

These cases must not be handled identically. A missing deprivation-band value
is an unknown attribute. A missing submission row is an event defined through
the relationship between attempts, assessments, and recorded submissions.

\verified{The saved audit contains 1,111 missing deprivation-band values and 11
assessments without due dates. The missing due dates are commonly associated
with exams, so those assessments are not treated as an eligible ``next
assessment'' when the later project defines a prediction outcome.}

\section{Representational limitations}

\begin{itemize}
  \item Module codes are anonymized, which limits subject-specific
    interpretation.
  \item OULAD represents selected Open University offerings, not all education.
  \item VLE records omit offline study and unrecorded activity.
  \item Final outcomes and unregistration dates do not contain every reason for
    an outcome.
  \item Demographic fields are sensitive contextual variables, not explanations
    of performance.
\end{itemize}

\section{Knowledge check}

\begin{enumerate}
  \item Why is \texttt{id\_student} not a student-attempt key?
  \item What does one row in \texttt{studentVle.csv} represent?
  \item How is an absent submission row different from a \texttt{NULL} score?
  \item Why should negative activity days not be deleted automatically?
  \item Which table connects a submission to its module-presentation?
\end{enumerate}

\begin{exercisebox}{Exercise 3.1: Build the grain and key worksheet}
Without looking back at \cref{tab:source-grains}, create a seven-row table with
these columns:
\begin{enumerate}
  \item source file;
  \item one-row meaning;
  \item candidate key;
  \item parent table;
  \item possible null or absent-row issue; and
  \item one analytical question the table can support.
\end{enumerate}
Then compare your table with the project source and correct it.
\solutionref{sol:source-grains}
\end{exercisebox}

\begin{exercisebox}{Exercise 3.2: Trace one learner attempt}
Choose the conceptual path from a student attempt to:
\begin{enumerate}
  \item its registration;
  \item its assessments;
  \item its recorded submissions; and
  \item its VLE interactions.
\end{enumerate}
List the columns required at each transition and state where row multiplication
could occur, meaning one starting row could turn into several matched rows. Do
not write the join yet.
\solutionref{sol:attempt-trace}
\end{exercisebox}

\begin{commonmistakes}
\begin{itemize}
  \item Calling 32,593 attempt rows ``32,593 students.'' The verified unique
    student count is 28,785.
  \item Treating \texttt{?} as literal text after staging.
  \item Assuming every assessment has a due day.
  \item Joining VLE resources by site ID without module-presentation context.
  \item Treating absence of a submission row as a database import failure.
\end{itemize}
\end{commonmistakes}

\transferheading
The seven-table worksheet generalizes directly to online-store events, health
claims, financial transactions, and app-usage logs. Before analysis, map every
source's row meaning, identifying columns, related parent record, time
representation, and way that values can be missing.

\begin{summarybox}
OULAD contains seven linked sources at different grains. Module, presentation,
student, assessment, site, and relative day combine differently across tables.
Student ID alone is not an attempt key. Missing cells and absent relationship
rows carry different meanings. Reliable analysis begins by defining one row,
its identifying columns, and its limitations.
\end{summarybox}
